Bett had carried water to three men in a cooper’s yard because nobody had told her not to, and because the alternative was sitting in her own room listening to the wall.

That was the whole of how she came to be there. There had been no order and no organizing. A boy had come up Cobbler’s Square at midday shouting that the healers had run out of hands, and Bett had gone down to the works with a bucket and a ladle because she was fifty-one years old and strong in the arms and had nothing else in the world anyone wanted, and she had spent the afternoon walking water along a line of straw pallets in a yard that smelled of oak shavings and blood.

The men were Northern. She could not understand more than half of what any of them said. One of them, a boy of about seventeen with something wrong with his hip, had asked her three separate times whether the wall was holding, and she had said yes three times, and she had no idea whether it was true and had decided that saying it was kinder than the alternative.

Then somebody outside the yard started screaming about \textit{behind}, and Bett turned round to see what the shouting was about, and the cellar hatch was open.

\sceneBreak

She knew what was in the cellars.

That was the thing. That was the thing she understood in the first half second, before the hand came up, before any of it — because Bett had spent the bad winter three years past on the corpse carts, and she had spent the summer before that on the pits after the Sept came down, and she was one of perhaps four hundred people still living in King’s Landing who knew precisely how many dead were under the ground she was standing on, and where, and how deep, because she had helped put them there with her own two hands.

They had not been graves. That was the part the septons had been careful with their words about. There had been no room for graves and no time and no wood for pyres, and so the city had done what the city always did when the dying outran the burying: they had opened the old cellars.

King’s Landing was built on King’s Landing. Everyone in Flea Bottom knew that. Under the streets there were older streets, and under those there were the cellars of houses that had burned in the reign of somebody nobody remembered, and under \textit{those} there were the pits, layer on layer, three hundred years of a city too crowded to bury anyone properly and too poor to burn them. When the winter had come and the bread had gone, the guildmen had come round with lanterns and picked out which cellars to open, and Bett had gone down the ladders with a scarf around her face and helped stack them, and had been paid in barley for it, and had eaten the barley.

She knew what was in the cellars. She had put a great deal of it there.

“Get up,” she said to the seventeen-year-old with the hip. “Get up now. Get \textit{up}.”

\sceneBreak

The first one out of the hatch was a man she did not know, and that was a mercy that lasted about forty seconds.

He came up slowly. That was what she would have told anyone about it afterward, if there had been any afterward: they came up \textit{slowly}, and it was worse than fast. A thing that came at you fast, you fought or you ran. This came up out of a cellar hatch in a cooper’s yard one arm at a time, patient as a man climbing out of a boat, and stood in the middle of the straw and the water buckets, and \textit{looked around}.

She had time to see the rope still round his ankles. They had tied ankles, that winter, to lower them down.

Then the second and third came up, and the fourth, and after the fourth they stopped coming up one at a time and started coming up the way water comes up.

She got the boy with the hip over her shoulder because there was nobody else to do it. She got the other one, the one who could walk, by the wrist. The third man on the straw was already dead and had been for an hour, and Bett looked at him for exactly as long as it took to be certain of it and then left him there, and would think about that later, in the last minutes she had, and would not be able to decide.

She went out of the yard into Cobbler’s Square, and the square told her the size of it.

\sceneBreak

It was not the works. It was not the wall or the trench or any of the places three days of shouting men had been about. It was the \textit{city}.

Every bell in King’s Landing was ringing and they were not ringing together. Bett came out of the cooper’s yard into a square she had crossed four times a day for thirty years, and there were people running in three directions and none of them were running from the north. A gap had opened in the paving outside the fishmonger’s where the old drain went down. The doors of the sept on the corner — small, poor, a sept for people who could not get into the big ones — stood wide open, and there was something in the doorway of it that was not a septon, and behind it there was a great deal of movement in the dark where the crypt stair went down.

She had been at a wedding in that sept. She had been at four funerals in it.

“North gate!” somebody was shouting. “Make for the north — “

“There’s nothing at the north gate,” Bett said, to nobody, because the man was already gone. “It’s not \textit{coming in}. You stupid — it’s already in. It’s been in the whole time.”

That was the thing nobody in the square had understood yet, and Bett had understood it in a cooper’s yard forty seconds after the hatch lifted, because she was the one who knew where the bodies were.

There was no direction to run.

\sceneBreak

She got them as far as Pisswater Bend before the boy with the hip started to go grey, and she put him down against a wall in a doorway she knew — Marra’s doorway, Marra who sold thread and had been dead since the spring — and she stood in the street and made herself think.

\textit{Up.}

That was the whole answer and it came to her clean. They came up out of cellars and drains and crypt stairs. They came up. They had come up slow, every one she had seen, and they had come up \textit{out of the ground}. So: up. Anywhere with a ladder and no cellar under it. The rope-walk at the end of Pisswater had a loft over it and one stair and no cellar at all because it had been built on the rubble of something that had burned.

“There,” she said. “That. Can you climb it? Say yes.”

“Aye,” the boy said, who could not have understood more than half of what she said, and did climb it, mostly, with Bett underneath him pushing.

She got eleven people into that loft in the next quarter hour. She would not have been able to say afterward how. Some of them she pulled off the street bodily. One was a woman with a baby who would not stop trying to go back for her sister, and Bett had hit her — hard, open-handed, across the face — and pushed her at the ladder, and the woman had gone up it weeping and calling her a bitch and Bett had thought \textit{good, be alive and call me a bitch} and gone back down for the next one.

The eleventh was a girl of about nine that she found sitting on a mounting block with her hands over her ears.

Then she went back out for a twelfth, and did not get one, because that was when she saw the small ones.

\sceneBreak

Nobody was watching the ground.

That was the thing. Every man on that street who had a weapon — and there were only three or four, a butcher with a cleaver, somebody’s grandfather with a spear from the last war — every one of them was standing with the point up and out at chest height, watching the street ahead for a man-sized shape, because that was what a person did, that was what fighting looked like. And the dead of the bad winter had been the old and the very young, because that was who a famine took first and always, and they were coming up the street at knee height.

They were fast. The grown ones came up slow and moved slow. The small ones did not. Bett never understood why and had about six minutes left to wonder about it.

She saw the butcher go down. He did not see the thing that took him at all; he was watching a doorway. He went down onto one knee with a noise she had never heard a man make and there were three of them on him before he got the cleaver round, and none of them came higher than his belt.

“\textit{Down!}” Bett was screaming. “\textit{Look down — look at the ground, look at the — }”

\sceneBreak

She found the child in Ragman’s Alley, between the tannery wall and the barrels, and she stopped running.

Bett had buried a son. His name was Tomm and he had been six and it had been the bad winter, the third month of it, and there had been no wood and no coin and no priest willing to come to Flea Bottom for free, and Bett had done what everyone in her street had done. She had carried him herself, because she would not let the cart take him, and she had gone down the ladder into the cellar under the old cooperage with him in her arms and she had put him somewhere near the wall where it was dry, and she had said the words she knew, which were not the right words, and she had climbed back up and gone to work.

She had known where he was for three years. She had walked over him four times a day for three years, going to the market.

He was standing in Ragman’s Alley with his arms at his sides.

Bett stopped, and she stood, and the whole enormous noise of the city fell away from her for a moment and she was aware of nothing at all except that he still had the shirt on. She had put him down in his shirt because it was all he had and it seemed wrong to take it. It was still on him. It was mostly still on him.

“Tomm,” she said.

She was not confused. She wanted that understood, in whatever part of herself was still keeping accounts. She was not one of those poor fools she had seen twice already in the last half hour, standing in the street with their arms out to something that used to be their wife. She knew exactly what was in that alley and she knew it was not her son and had not been her son for three years, and she knew, with the entire practical intelligence that had got fifty-one years out of a city that had never once given her anything, that the correct thing to do was turn around and go back up the ladder into the loft with the eleven people in it.

She was not confused at all. She simply found that she was not going to do that.

“You come here, then,” she said. “Come here to me.”

\sceneBreak

It was very fast.

She got her arms round him. She had thought about this in the two seconds she had to think in — that if she got her arms round him and got him against the tannery wall, then the eleven people in the loft at the end of Pisswater would have another minute, and possibly two.

He was so light. That was what she had time to notice. She had carried him down a ladder once and he had been heavier then than he was now.

It took considerably longer than she expected, and hurt in a place she had not anticipated, and at some point she was on the ground in the wet between the barrels and the wall and the noise of the city had come back very loud, and she was thinking, in a distant and unbothered way, about barley. About being paid in barley to stack the cellar under the cooperage. About how she had eaten it.

There were two of the small ones on her now and she could not feel either of them properly.

\textit{The loft,} she thought. \textit{Eleven. Twelve with the girl. Twelve.}

She could hear the bells, still not ringing together, all across the city she had lived in her whole life and never once left.

And then, right at the end, with the cold coming up through her the way water came up through the paving, Bett understood one more thing — one last piece of arithmetic, arriving too late to be any use to anyone, the way the useful things always arrived.

\textit{I’ll be getting up.}

She had helped stack four cellars. She knew exactly how the sums went.

\textit{Oh,} she thought. \textit{Oh, that’s not fair. That’s not — }